%!TEX root = ./../../springer-mono-format/book.tex
%!TEX root = ./../../thesis-trier.tex
  
  \todo[inline]{viel auskommentiert}
  
 Here, we discuss some other complexity results on problems related to synchronization from the literature.
 In view of the correspondence between automata
 and transformation semigroups from \cref{prelim:remarks_trans_semigrp_semiautomata}, 
 we also discuss related problems on transformation semigroups. 
 For notions from (parameterized) complexity
 theory or automata theory
 not defined before, we refer to the cited papers
 or other literature as, for example,~\cite{DBLP:books/daglib/0023084}
 for classical computational complexity theory,
 \cite{FluGro2006} for parameterized complexity theory, or~\cite{DBLP:reference/hfl/Pin97} for automata theory.
  
 In this section, we do not discuss complexity issues
 related to the constrained synchronization problem
 introduced in \cref{overview:sec:csp} as this is done in \cref{part:constrsync} of the present work.
  
\paragraph{Further Complexity Theoretical Results.} \NP-completeness of \DFASWS\ 
was established for the following classes of input automata in the literature:
\begin{enumerate}
\item \label{prelim:compl:discussion:pto} Binary partially ordered automata with a unique trap state~\cite{Eppstein90,Rys80};

\item SI-automata over an unbounded alphabet~\cite{DBLP:conf/fun/FernauH022,Mar2009b};

\item Partially ordered commutative automata~\cite{DBLP:conf/fun/FernauH022,DBLP:journals/jcss/FernauHV15,Mar2009b};

\item Eulerian\footnote{An automaton $\semiautomat{A}{Q}{\Sigma}{\delta}$ is \emph{Eulerian}
	if $\sum_{x\in \Sigma} | \{ p \in Q \mid \applytrans{\delta}{p}{x} = q \} |
	= |\Sigma|$ for each $q \in Q$.}
automata over a binary alphabet~\cite{DBLP:journals/iandc/Vorel17};

\item  Circular automata over a binary alphabet~\cite{DBLP:journals/tcs/Martyugin12};

\item Automata whose automaton multigraph is a two-terminal series-parallel multigraph with loops (TTSPL automata) over an alphabet with at least four letters~\cite{DBLP:journals/ita/BruchertseiferF21}.
\end{enumerate}

Rystsov~\cite{Rys80} and Eppstein~\cite{Eppstein90}, independently,
gave the same reduction to show \NP-hardness by constructing
a binary partially ordered automaton.
Essentially the same reduction was also given by Salomaa~\cite{DBLP:journals/tcs/Salomaa03}
to show \NP-hardness of the problem whether a specific constant function can be composed out of a given
set of functions.
The reduction also appears as a special case in a paper 
by Goral\v{c}\'yk and Koubek~\cite{GorKou95} where 
the question of realizing idempotent transformations
of a certain rank was investigated.
Martyugin~\cite[Prop.~4.2]{Mar2009b} noted that the
automaton resulting from this reduction has a trap state,
is aperiodic, has a $\mathcal J$-trivial transition monoid,
and belongs to the class of automata preserving a linear order on the states
except at a single trap state, thus giving \NP-hardness for these classes as well.\footnote{He called the latter automata partially monotonic automata, but since these automata
	have a totally defined transition function
	we avoid this term here
	to not cause confusion with partial automata where the transition function
	can be a partial function introduced later in \cref{part:constrsync}.}
An automaton $\semiautomat{A}{Q}{\Sigma}{\delta}$ is called \emph{locally confluent}
if, for any $q \in Q$ and $a,b \in \Sigma$,
there exists $u \in \{a,b\}^*$ such that $\delta(q, au) = \delta(q, bu)$.
Clearly, every synchronizing automaton
over a binary alphabet is locally confluent.
So, by the Rystsov-Eppstein-Salomaa reduction,
\DFASWS\ is \NP-complete for locally confluent partially ordered automata.
The \NP-hardness result
on TTSPL automata
was shown in~\cite{DBLP:journals/ita/BruchertseiferF21}
by adding two additional letters and additional states to obtain a TTSPL automaton
from the Rystsov-Eppstein-Salomaa construction.

Eppstein~\cite{Eppstein90} showed that for oriented $n$-state automata\footnote{This is the terminology due to Ananichev and 
	Volkov~\cite{DBLP:journals/tcs/AnanichevV04} which we use here as well, Eppstein himself called these automata monotonic; see 
	also the discussion on page~\pageref{discussion:monotonic}.}  
over the alphabet~$\Sigma$ the problem \DFASWS\ is solvable
in time $O(|\Sigma| n^2 )$. Every monotonic automaton is oriented~\cite{DBLP:journals/tcs/AnanichevV04},
yielding polynomial time solvability for monotonic automata as well.
Martyugin~\cite{Mar2009b} showed that \DFASWS\ is solvable in $O(n^{|\Sigma|} \log(n))$
for commutative automata and in $O(n)$ for binary SI-automata with~$n$ states (by checking
whether it equals one of the two SI-automata from \cref{prop:structure_simp_idem_binary}).

Note that, while \DFASWS
remains \NP-complete even for partially ordered commutative input automata (\cref{compl:thm:comm-syncshort}),
as we have shown in \cref{compl:sec:subclasses}, the complexity of \DFAS 
drops from \NL-completeness in the general case
to~$\AC^0$ on partially ordered input automata (\cref{compl:thm:pto_sync_in_AC0})
and to~$\LOGSPACE$-completeness for solvable input automata (\cref{thm:SYNC-solv}).

Olschweski and Ummels~\cite{DBLP:conf/mfcs/OlschewskiU10}
proved that deciding whether a given automaton
has a shortest synchronizing word of length \emph{exactly} $k$
is \DP-complete. Furthermore, they proved
that computing the length of a shortest synchronizing word,
a synchronizing word of length $k$,
and a shortest synchronizing word
is~$\FP^{\NP[\log]}$-complete, \FNP-complete, and $\FP^{\NP[\log]}$-hard and in $\FP^{\NP}$, respectively.
With respect to approximability, for $n$-state automata,
Gawrychowski and Straszak~\cite{DBLP:conf/mfcs/GawrychowskiS15}
proved that for every $\epsilon > 0$
it is \NP-hard to approximate the length of a shortest
synchronizing word within a factor of~$n^{1 - \epsilon}$ even for a fixed binary alphabet.
This result is optimal, as the Eppstein-Natarajan algorithm~\cite{Eppstein90,DBLP:conf/focs/Natarajan86}
mentioned in \cref{prelim:sync:algorithmic} 
gives an $(n-1)$-approximation algorithm~\cite{DBLP:conf/wia/GerbushH10}.
More generally, Gerbush and Heeringa~\cite{DBLP:conf/wia/GerbushH10} 
gave a $O(k|\Sigma|n^k + n^4/k)$-time $\lceil \frac{n-1}{k-1} \rceil$-approximation algorithm, where $k \ge 2$ 
denotes the size of the intermediate sets used for compression.
Historically, before the optimal result by Gawrychowski and Straszak~\cite{DBLP:conf/mfcs/GawrychowskiS15},
under the assumption $\PTIME \ne \NP$, 
Berlinkov firstly proved 
that there exists no polynomial time~$O(1)$-approximation~\cite{DBLP:conf/csr/Berlinkov10,DBLP:journals/mst/Berlinkov14}
over a fixed binary alphabet.
Gerbush and Heeringa~\cite{DBLP:conf/wia/GerbushH10} 
then proved that there is no polynomial time $O(\log n)$-approximation
over an unbounded alphabet, which
was improved by Berlinkov~\cite{DBLP:conf/dlt/Berlinkov14} 
to a fixed binary alphabet.

Fernau and Hoffmann~\cite{FernauHoffmann19} introduced, for a given partial order $\prec$ on $\Sigma^*$,
the following \emph{extension variant}
of the synchronization problem.

\decisionproblem{{\sc ExtDFA-SW-$\prec$} (Fernau and Hoffmann~\cite{FernauHoffmann19})}
{Automaton $\semiautomat{A}{Q}{\Sigma}{\delta}$, $u \in \Sigma^*$.}
{Is there $w \in \Syn(\mathcal A)$ minimal w.r.t.~$\prec$ in $\Syn(\mathcal A)$ such that $u \prec w$?} 

For the prefix/suffix/lexicographic orders the problem is in \PTIME, 
for the scattered subsequence order even on commutative automata (and even for TTSPL automata)
it is \NP-hard.
For the parameter $|u|$ and the scattered subsequence order
it is $\W[3]$-hard even on commutative automata (and even for TTSPL automata).
The complementary problem that switches {\rYES} and \rNO answers
for the length-lexicographic order is \NP-complete even for TTSPL automata
and in $\WNL\cap \W[P]\cap \textsf{A}[2]$ w.r.t.~$|u|$.
For more information, we refer to~\cite{DBLP:journals/ita/BruchertseiferF21,DBLP:journals/mscs/FernauB22,FernauHoffmann19}.

Maslennikova~\cite{DBLP:journals/ijfcs/Maslennikova19}
showed that it is \PSPACE-complete
to decide whether a given regular ideal language
is the set of synchronizing words of an automaton
with at most four states over a binary alphabet.

Türker and Yenigün~\cite{DBLP:journals/ijfcs/TurkerY15}
introduced and investigated
variants of the synchronization problem, 
the following problems being their main problems.

\decisionproblem{{\sc MinimumSynchronizableSubautomaton} (Türker and Yenigün~\cite{DBLP:journals/ijfcs/TurkerY15})}
{Automaton $\semiautomat{A}{Q}{\Sigma}{\delta}$ and $k \ge 0$.}
{Does there exist $\Gamma \subseteq \Sigma$, $|\Gamma| \le k$, such that $\mathcal A_{|\Gamma}$ is synchronizing?}

\decisionproblem{{\sc ExclusiveSyncWord} (Türker and Yenigün~\cite{DBLP:journals/ijfcs/TurkerY15})}
{Automaton $\semiautomat{A}{Q}{\Sigma}{\delta}$, $S, T \subseteq Q$.}
{Does there exist $u \in \Sigma^*$ such that $\applytrans{\delta^{-1}}{T}{u} = S$?}

\decisionproblem{{\sc MaxF-ReducingWord} (Türker and Yenigün~\cite{DBLP:journals/ijfcs/TurkerY15})}
{Automaton $\semiautomat{A}{Q}{\Sigma}{\delta}$, $F \subseteq Q$, $k \ge 0$.}
{Does there exist $S \subseteq Q$, $|S| \ge k$, $u \in \Sigma^*$ such that $\applytrans{\delta}{S}{u} \subseteq F$?}


They proved that {\sc MinimumSynchronizableSubautomaton}
is \NP-complete\footnote{We omitted a cost function on the letters 
	mentioned in~\cite{DBLP:journals/ijfcs/TurkerY15}, which does not change \NP-completeness.
Also, in~\cite{DBLP:journals/ijfcs/TurkerY15} only \NP-hardness is stated,
but, as  noted in~\cite{DBLP:journals/ita/BruchertseiferF21},
containment in \NP\  follows as we can guess $\Gamma$ and check synchronizability of $\mathcal A_{|\Gamma}$ in polynomial time.}
%by guessing $\Gamma$ and checking synchronizability of $\mathcal A_{|\Gamma}$ in polynomial time.} 
(which holds even for TTSPL automata as shown in~\cite[Cor.~5.9]{DBLP:journals/ita/BruchertseiferF21}), 
{\sc MaxF-ReducingWord} is \NP-hard
for monotonic automata and
{\sc ExclusiveSyncWord}
is \PSPACE-complete. The reduction for {\sc MaxF-ReducingWord}
was challenged by Ryzhikov~\cite[p.~80]{Ryzhikov19a} 
as ``they reduce the {\sc $N$-Queens Puzzle}~\cite{DBLP:journals/dm/BellS09} 
to this problem to prove its \NP-hardness. However, their poof
is unclear, since the input has size $O(\log N)$,
and the output size is polynomial in $N$.``
Ryzhikov~\cite[Prop.~3]{Ryzhikov19a} showed
that {\sc MaxF-ReducingWord} is \PSPACE-complete for binary automata.
In the context of parameterized complexity, 
{\sc MinimumSynchronizableSubautomaton}
was shown to be $\W[2]$-hard even for TTSPL automata~\cite[Cor.~5.7 \& 5.9]{DBLP:journals/ita/BruchertseiferF21} 
and in $\WNL \cap \W[P]$~\cite[Thm.~12]{DBLP:journals/mscs/FernauB22}.

\Cref{compl:sec:subclasses:known-results}
collected known results on the subset mapping problems.
Next, we give some further background.	
\SIS, \SFS, and
\ST were introduced, and shown to be \PSPACE-complete,
by Rystsov~\cite{Rys83} under the names 
\emph{(global) inclusion problem 
	for non-initial  au\-tom\-a\-ta} (\emph{NIA})~\cite[Thm.~2.1]{Rys83},
\emph{(global) inclusion problem for weakly-initial automata with a single final state} (\emph{$\textsc{WIA}_1$})~\cite[Cor.~3.3]{Rys83},
and \emph{(global) inclusion problem for weakly-initial automata} (\emph{WIA})~\cite[Thm.~2.1]{Rys83}, 
%\begin{enumerate}
%\item \emph{(global) inclusion problem 
%	for non-initial  au\-tom\-a\-ta} (\emph{NIA})~\cite[Thm.~2.1]{Rys83},
%
%\item \emph{(global) inclusion problem for weakly-initial automata with a single final state} (\emph{$\textsc{WIA}_1$})~\cite[Cor.~3.3]{Rys83},
%
%\item \emph{(global) inclusion problem for weakly-initial automata} (\emph{WIA})~\cite[Thm.~2.1]{Rys83}, 
%\end{enumerate}
respectively.
Independently, Luks and McKenzie~\cite{LuksM88}
later introduced \ST for permutation groups that asks, %\todo{zusammenhang mit automaten hier schon erwähnen, oder das oben in T_A kapitel alles erwähnen?}
for a permutation group $G$ on~$[n]$ given by generators and subsets~$A,B \subseteq [n]$,
whether  $\applyfunc{A}{g} = B$ for some $g \in G$. 
They proved that the problem is in \NC\ for  solvable permutation groups. % (see~\cite{LuksM88} for the definition of solvable groups).
For groups, the sets $A$ and $B$ must have the same size
and $\applyfunc{A}{g} = B$ is equivalent to $g(A) \subseteq B$.
We note that the automata-theoretical analogous problem 
with equality instead of set inclusion
%where set inclusion is replaced by equality
was investigated by Rystsov under the name \emph{reachability problem for weakly-initial automata}
and proved to be \PSPACE-complete~\cite[Thm.~3.2]{Rys83},
even when the starting set is the whole state set and the target has size two.\footnote{The case
that the starting set is the whole state set was also
shown to be \PSPACE-complete independently by Bondar and Volkov~\cite[Thm.~3]{BondarV16}.}
In~\cite{GOLASZEWSKI1989431}, reachability problems in certain discrete event systems
were investigated that translate, in our language, to the problem \DFAINT 
for unary permutation automata and with different types of final state sets,
giving polynomial-time solvable and \NP-complete problems.


Blondin et al.~\cite{BlondinKM16} investigated different variants and generalizations
of \DFAINT.
They named \ST the problem whether there exists, in a transformation monoid given by a list of generators,
a transformation mapping a given set into a set included in another given set. 
\ST as defined
here is the automata-theoretical version of this problem, i.e., we opted to use
the naming suggested by the work of Blondin et al.~\cite{BlondinKM16}.
%Vorel~\cite{DBLP:journals/ijfcs/Vorel16} has shown that \SFS 
%is \PSPACE-complete even on strongly connected binary automata. \todo{noch genauer bezüglich alphabetgrößen besprechen?}
For a fixed number $k \ge 1$ of automata, \DFAINT 
is \NL-complete~\cite[p.~464]{HolzerK11}.


\IN{extending problem}
\IN{totally extending problem}
\IN{resizing problem}
\IN{avoiding word}
Berlinkov, Ferens and Szyku\l a~\cite{BerlinkovFS21}
classified and investigated various further subset mapping problems.
For $\semiautomat{A}{Q}{\Sigma}{\delta}$ and $S \subseteq Q$,
their problems fall into the following broad main categories:
\emph{extending problems}, \emph{totally extending problems}, and \emph{resizing problems}
that ask for $w \in \Sigma^*$
such that $|\applytrans{\delta^{-1}}{S}{w}| > |S|$,
$\applytrans{\delta^{-1}}{S}{w} = Q$, and
$|\applytrans{\delta^{-1}}{S}{w}| \ne |S|$, respectively.
They investigated those problems with and without a binary encoded bound on~$|w|$
and, for the former two categories,
with and without a fixed parameter $k$ such that  $|S| \le k$ or $|S| \ge n - k$, which is used,
for example, in runtime bounds of algorithms.
%algorithms were given whose runtime 
%is parameterized by a fixed $k$ for the case $|S| \le k$
%and for the case $|S| \ge n - k$.
A word $w \in \Sigma^*$ \emph{avoids} $S \subseteq Q$ 
if $\applytrans{\delta}{Q}{w} \cap S = \emptyset$.
As this is equivalent to $\applytrans{\delta^{-1}}{Q \setminus S}{w} = Q$,
totally extending problems can be rephrased
as problems that ask for avoiding words of the complement.
Furthermore, 
the authors considered \emph{compression}, i.e., $|\applytrans{\delta}{S}{w}| < |S|$,
which they showed to be polynomial-time solvable, and \emph{shrinking}, i.e., $|\applytrans{\delta^{-1}}{S}{w}| < |S|$,
which is equivalent to extending~$Q \setminus S$.
As $\applytrans{\delta^{-1}}{S}{w} = Q$
is equivalent to $\applytrans{\delta}{Q}{w} \subseteq S$,
the problem of total extension without parameters is
equivalent to \SIS, which was investigated earlier by Rystsov~\cite{Rys83}. %(the authors of~\cite{BerlinkovFS21} do not seem to mention this connection).
%While citing~\cite{Rys83} the authors of~\cite{BerlinkovFS21} do not seem to mention this connection.
For a full overview of the results we refer to their work~\cite{BerlinkovFS21}.

\IN{synchronization diversity $k$}
A recent trend in fixed parameter tractability~\cite{DBLP:journals/ai/BasteFJMOPR22}
and the field of artificial intelligence~\cite{DBLP:conf/aaai/HebrardHOW05,DBLP:conf/aaai/IngmarBST20}
asks for \emph{diverse solution sets}, giving a user more freedom of selection of a particular
solution to fit his/her needs.  Let $\mathcal A$ be an automaton.
Arrighi et al.~\cite{DBLP:conf/aaai/ArrighiFO023}
extended automata synchronization to the diversity setting \todo[inline]{prüfen.}
by saying a finite set $W \subseteq \Syn(\mathcal A)$
has \emph{synchronization diversity} $k$
if two distinct words from $W$ have edit distance at least $k$ and each word in $W$
is minimal with respect to the scattered subsequence order in $\Syn(\mathcal A)$.
Then, they proved that deciding whether an automaton
admits a sets of synchronization diversity at least $k$
is co-\NP-hard and that computing
the diversity of the set of words minimal w.r.t.~scattered subsequence order
in $\Syn(\mathcal A)$
is hard for the counting complexity class $\# \PTIME$.

\paragraph{Parameterized Complexity.} 
 The following results on the parameterized complexity of \DFASWS are known (the paramterized complexity class
 \WNL is a lesser known
 class defined via the number of tape cells used
 by a nondeterministic Turing machine~\cite{DBLP:journals/disopt/Guillemot11}):
 
 \begin{enumerate}
 \item \DFASWS is in \FPT\ w.r.t.~$(|\Sigma|, l)$~\cite[Thm.~7]{DBLP:journals/jcss/FernauHV15},
 as we can try all $|\Sigma|^l$ words of length~$l$, 
 and the number of states~$n$~\cite[Thm.~10]{DBLP:journals/jcss/FernauHV15},
 as we can search in the power automaton with~$2^{n}$ states and~$|\Sigma|2^{n}$ transitions.
 Furthermore, unless the polynomial-time hierarchy collapses on the third level,
 \DFASWS admits no polynomial-size kernel
 when parameterized by $(|\Sigma|, l)$~\cite[Cor.~4.6]{DBLP:journals/ita/BruchertseiferF21}
 or by~$n$~\cite{DBLP:journals/dmtcs/VorelR15}.
 	
 \item Bruchertseifer and Fernau~\cite{DBLP:journals/mscs/FernauB22}
 introduced the parameterized complexity class \WSYNC 
 as the class of problems fpt-reducible
 to \DFASWS with parameter being the length-bound~$l$. 
 They showed that \WSYNC 
 sits between~$\W[2]$ and $\WNL \cap \W[\textsf{P}] \cap \textsf{A}[2]$,
 see~\cite[Thm.~3]{DBLP:journals/jcss/FernauHV15},
 \cite[Thm.~4.2 \& Prop.~4.4]{DBLP:journals/ita/BruchertseiferF21},
 and~\cite[Thm.~3,4 \&~5]{DBLP:journals/mscs/FernauB22}. 
 \end{enumerate}
 Montoya and Nolasco~\cite{MontoyaN17,DBLP:conf/lata/MontoyaN18} 
 proved that \SFS with a length bound on the solution word is \WNL-complete
 when parameterized by the size of the input set.

Wareham~\cite{DBLP:conf/wia/Wareham00} investigated \emph{bounded automata intersection}
with a parameter~$l$ asking for a solution word
of length exactly $l$
accepted by $m$ automata
each with at most~$n$ states.
He obtained containment in \FPT\  for the following parameter combinations:
\begin{enumerate}
	\item \label{prelim:classes:Wareham:1} $(|\Sigma|, l)$ by trying all words in $\Sigma^{l}$,
	
	% \item $(m,n,l)$ via the product automaton of size $n^m(l+1)$
	% of all input automata and an automaton accepting $\Sigma^{\le l}$,
	
	\item \label{prelim:classes:Wareham:2} $(m,n)$ via the product automaton of size $n^m$
	of all input automata, 
	
	\item  \label{prelim:classes:Wareham:3} $(n, |\Sigma|)$ by noting that there are at most $n^{|\Sigma|n}$
	different automata, so that we can assume $m \le n^{|\Sigma|n}$, and use
	the reasoning as for $(m,n)$.
\end{enumerate}
In fact, Wareham stated the above
results under the assumption $l \le n$. 
Clearly, this restriction is not necessary for~\ref{prelim:classes:Wareham:1} to hold.
But it is also not necessary for~\ref{prelim:classes:Wareham:2},
and so~\ref{prelim:classes:Wareham:3}, 
since the $l$-th power of the transition matrix of the automaton graph of the product
automaton mentioned in~\ref{prelim:classes:Wareham:2} 
can be computed with $\log(l)$ multiplications
by a square-and-multiply algorithm,
and this number is polynomial even for a binary encoding of~$l$.
For hardness, he obtained the following results:
$W[1]$-hard for $(m,l)$, $W[2]$-hard for $(n, l)$,
$W[t]$-hard for all $t \ge 1$ for the sole parameter $m$,
and not in~$\XP$ for the sole parameter~$|\Sigma|$ (under the assumption $\PTIME \ne \NP$);
the construction also implies this for the sole parameter $n$, as noted in~\cite{DBLP:journals/corr/abs-2108-05244}.
Further, Bruchertseifer and Fernau~\cite[Thm.~9]{DBLP:journals/mscs/FernauB22}, Fernau et al.~\cite[Cor.~2.2]{DBLP:journals/corr/abs-2108-05244}, and Guillemot~\cite[Prop.~10]{DBLP:journals/disopt/Guillemot11} 
proved \WSYNC-completeness w.r.t.~$l$, $W[1]$-completeness w.r.t.~$(m, k)$, and \WNL-completeness w.r.t.~$m$, respectively.
Moreover, Fernau et al.~\cite[Thm.~5.8]{DBLP:journals/corr/abs-2108-05244} 
proved $W[1]$-completeness w.r.t.~$m$
for automata accepting sparse languages even for
unary automata (which accept sparse languages).\footnote{See \cref{constr_sync:sec:growth} for a definition of sparse languages.}
We remark that the co-$W[2]$-hardness proof in~\cite[Prop.~6.3]{DBLP:journals/corr/abs-2108-05244}
is wrong.
For further information on synchronization
problems, parameterized complexity,
and also runtime lower bounds derived from ETH and SETH 
we refer to~\cite{DBLP:journals/ita/BruchertseiferF21,DBLP:journals/jcss/FernauHV15}.

Later, in \cref{cor:aut_for_sync_words,compl:cor:intersection_waa_com},
we will show
that \DFAINT and \DFASWS 
are solvable in time $O(m|\Sigma|n^{|\Sigma|+1})$ on  $m$ commutative partially ordered input automata with at most $n$ states
and in time $O(n^{|\Sigma|+2})$
on commutative $n$-state automata, respectively.\footnote{Martyugin~\cite{Mar2009b} obtained an algorithm running in time 
	$O(|\Sigma|n^{|\Sigma|}\log(n))$.}
So, both problems are in the parameterized complexity class \XP\ w.r.t.~$|\Sigma|$
for the mentioned types of input automata.
The \XP-result for \DFAINT on commutative partially ordered automata
has been extended in~\cite{Arrighi2021}
to  \emph{totally star-free nondeterministic automata} 
accepting commutative languages.


Lastly, testing whether $m$ binary automata accept a common word
is equivalent to the simulation of a nondeterministic Turing machine
that uses $m \log(n)$ bits of space~\cite{DBLP:journals/mst/KasaiI85} (see also~\cite{Galil76,DBLP:journals/tcs/KarakostasLV03,DBLP:conf/mfcs/LangeR92,Wehar2014hardness,Wehar2016}).
In~\cite{DBLP:journals/tcs/KarakostasLV03,Wehar2014hardness}, for each fixed~$m$
the resulting families of automata intersection problems have been used to obtain conditional separation results between $\NL$ and~$\PTIME$.	


\paragraph{Practical Implementations to Compute Shortest Synchronizing Words.}
As \DFASWS is \NP-complete
there exists no polynomial time algorithm computing a shortest synchronizing word
under the assumption $\NP \ne \PTIME$.
%we cannot expect to find an algorithm computing a shortest synchronizing word 
%that always runs efficiently.
However, people have tried to come up with various heuristics to obtain
algorithms that run efficiently
on most instances in practice.
The earliest implementations of exact
algorithms using several approaches were given
by Rho et al.~\cite{DBLP:conf/dac/RhoSP93},
Pixley et al.~\cite{DBLP:journals/tcad/PixleyJH94},
and Deshmukh and Hawat~\cite{Deshmukh94}.
Kisielewicz et al.~\cite{KisKowSzy2015}
have given an algorithm that combines multiple techniques and performs
well on most instances.
This algorithm
has been recently improved by Szyku\l a and Zyzik~\cite{DBLP:conf/esa/SzykulaZ22} (see also the related forward and backward scheme
by Szyku\l a and Roman~\cite{DBLP:journals/eswa/RomanS15}).
Algorithms
based on machine learning were proposed
by Podolak et al.~\cite{DBLP:conf/icaisc/PodolakRJ12,DBLP:journals/eswa/PodolakRSZ18},
and an evolutionary algorithm by Kowalksi and Roman~\cite{DBLP:conf/evoW/KowalskiR17}.
Tas et al.~\cite{DBLP:journals/isci/TasKY21} studied
the performance of some algorithms exploiting GPUs and parallelism on automata with up to a billion states.
Furthermore, in~\cite{DBLP:conf/modelsward/CirisciSSKKY18,DBLP:journals/eswa/KarahodaKY18,DBLP:journals/eswa/SaracAAKKY21}
the authors optimize existing
algorithms for real-life applications, for example in model-based testing.

The Natarajan-Eppstein algorithm~\cite{Eppstein90,DBLP:conf/focs/Natarajan86} mentioned in \SEC~\ref{prelim:sync:algorithmic}
runs in polynomial time and 
outputs some synchronizing word
that is guaranteed to have at most cubic length.
Two other algorithms with these properties
were suggested by Trahtman~\cite{Trahtman04}\footnote{In other works spelled as Trakhtman.}
and Roman~\cite{DBLP:journals/amc/Roman09}.
Kud\l acik et al~\cite{KUDLACIK201211746}
compare different exact and non-exact algorithms.

For more information, we refer to the introductory section of~\cite{KisKowSzy2015}
that contains further information and additional references
to implementations and algorithms.

\IN{{\sc TransformationRank}}
\IN{{\sc IdempotentRank}}
\IN{{\sc Membership}}
\paragraph{Related Problems on Transformation Semigroups.} 
Goral\v{c}\'{i}k and Koubek~\cite{GorKou95}
investigated the problem {\sc TransformationRank} (resp.~{\sc IdempotentRank}) 
which asks, for a set $T \subseteq \mathcal T_n$ and $r > 0$,
whether $\langle T \rangle_S$ contains a transformation (resp.~idempotent transformation)
of rank~$r$. For the fixed parameter $r = 1$, both problems
reduce to \DFAS.
Further, for each fixed $r > 0$ they proved that {\sc TransformationRank} is \NP-hard
and {\sc IdempotentRank} is solvable in time $O((r+1)n^{(2r+1)}|T|)$.
In general, {\sc IdempotentRank} was proved to be \PSPACE-complete.
The precise complexity of {\sc TransformationRank} is an \emph{open problem}.
A variant of {\sc IdempotentRank} with a 
length-bound for a product of transformations from~$T$
was proved to be \NP-complete for each fixed $r$.\footnote{Their proof~\cite[Thm.~5]{GorKou95} reduces 
	the problem to the special case $r = 1$
	and so, basically, \DFASWS. 
	They give a reduction from \THREESAT
	which is a little different than
	the Rystsov-Eppstein construction~\cite{Eppstein90,Rys80} 
	but similar to constructions
	from~\cite{DBLP:conf/fun/FernauH022,DBLP:journals/jcss/FernauHV15,FernauHoffmann19,Mar2009b}.}
In the realm of parameterized complexity theory,
Bruchertseifer and Fernau~\cite[Thm.~8]{DBLP:journals/mscs/FernauB22}
proved {\sc TransformationRank} complete for~\WSYNC.

Kozen~\cite{Kozen77} introduced 
{\sc Membership} that asks, for $T \subseteq \mathcal T_n$ and $t \in \mathcal T_n$,
whether~$t \in \langle T \rangle_S$.
Now, we obtain various variants of {\sc Membership}
by restricting~$T$ to generate monoids from some given class. 
Similarly, for the problem \DFAINT
by focusing the attention to automata whose transformation
monoids come from some common class of monoids,
we obtain variants of this problem
With respect to these variants,
Blondin et al.~\cite{BlondinKM16} observed the following relation.
\begin{quote}
	``[{\sc DFA-Intersection}] is equivalent to [{\sc Membership}] when both are intractable, but appears harder than [{\sc Membership}] when [{\sc Membership}] is efficiently solvable.''~\cite{BlondinKM16}	
\end{quote}
In fact, the complexity gap can be quite large.
For example, by results of Arvind and Vijayaraghavan~\cite{DBLP:conf/coco/ArvindV04} 
it follows that {\sc Membership} on abelian permutation groups 
is in $\NC^2$, whereas \DFAINT even on unary permutation
automata is \NP-complete~\cite{Beaudry88thesis,DBLP:journals/corr/abs-2108-05244,Galil76,MartensN07,Stockmeyer73} (see also \cref{compl:thm:comm_ST}).
More specifically, we have $\text{\sc Membership} \leACZEROm \DFAINT$ as for 
$t, t_1, \ldots, t_k \colon [n] \to [n]$
with $\semiautomat{A}{[n]}{\Sigma}{\delta}$
given by $\Sigma = \{t,t_1, \ldots, t_k\}$
and $\applytrans{\delta}{i}{x} = \applyfunc{i}{x}$ for $x \in \Sigma$
we have
\[
t \in \langle \{ t_1, \ldots, t_k \}\rangle_S
\Leftrightarrow 
\bigcap_{i=0}^{n-1} L(\mathcal A_{i, \{ \applyfunc{i}{t} \}}) \ne \emptyset.
\]

Now, suppose $\automat{A_i}{Q_i}{\Sigma}{\delta_i}{q_i}{F_i}$
with $i \in \{1,\ldots,k\}$ form an input to \DFAINT.
We can assume the state sets $Q_i$ are pairwise disjoint.
Let $t_a \colon Q_1 \cup \cdots \cup Q_k \to Q_1 \cup \cdots \cup Q_k$ for $a \in \Sigma$
be the transformations with $\applyfunc{q}{t_a} = \applytrans{\delta_i}{q}{a}$
for the unique $i \in \{1,\ldots,k\}$ such that $q \in Q_i$.
Then
\[
\bigcap_{i=1}^{k} L(\mathcal A_i) \ne \emptyset 
\Leftrightarrow 
\exists t \in \langle \{ t_a \mid a \in \Sigma \} \rangle_S \  \forall i \in \{1,\ldots,k\} \colon \applyfunc{q_i}{t} \in F_i. 
\]
So, we can solve $\bigcap_{i=0}^{n-1} L(\mathcal A_i) \ne \emptyset$ by guessing a transformation $t$
and then checking
checking whether $t \in \langle \{ t_a \mid a \in \Sigma \} \rangle_S$ and
$\applyfunc{q_i}{t} \in F_i$ 
for all $i \in \{1,\ldots,k\}$.
This reduction up to some ``nondeterminism''
makes the ``equivalent when intractable``-statement more precise.\footnote{The construction is due to Lukas Fleischer (personal communication).}
Note that, for classes generated by the $t_a$ such that {\sc Membership} is solvable in polynomial time,
this procedure gives an \FPT-algorithm for the parameter $\max\{ |F_i| \mid i \in \{1,\ldots,k\}\}$.

Closer inspection of Kozen's proof~\cite[Thm.~3.2.6]{Kozen77} reveals that
{\sc Membership} is \PSPACE-complete even 
for aperiodic transformation monoids as input.
For aperiodic transformation monoids, the {\sc Membership} problem
was classified into five classes giving problems that are in $\AC^0$,
$\PTIME$-~,~\NP-, or \PSPACE-complete, or \NP-hard, by Beaudry et al.~\cite{DBLP:journals/jacm/BeaudryMT92},
where \NP-hardness holds for aperiodic transformation monoids
of threshold at least two. Partial results
for general threshold-one monoids were obtained by Beaudry~\cite{DBLP:journals/iandc/Beaudry94}.
For commutative transformation monoids, 
Beaudry~\cite{DBLP:journals/iandc/Beaudry88} proved \NP-completeness.

For permutation groups, {\sc Membership}
was shown to belong to \PTIME by Furst et al.~\cite{DBLP:conf/focs/FurstHL80} (analyzing an algorithm given by Sims~\cite{Sims70}),
then to $\NC^3$ for abelian groups by McKenzie and Cook~\cite{DBLP:journals/siamcomp/McKenzieC87} 
and Mulmuley~\cite{DBLP:journals/combinatorica/Mulmuley87} 
(improved 
by Arvind and Vijayaraghavan~\cite{DBLP:conf/coco/ArvindV04} 
to containment in~$\textsf{ZPL}^{\textsf{ModL}}$,
which is in $\NC^2$).
Luks and McKenzie~\cite{DBLP:journals/jcss/LuksM88} proved membership in \NC for solvable groups, Luks~\cite{DBLP:conf/focs/Luks86} extended this to groups with bounded non-abelian composition factors, and Babai et al.~\cite{DBLP:conf/stoc/BabaiLS87} finally established membership in \NC for all groups.

Now consider the variant of {\sc Membership} parameterized by an integer $l$, 
where one asks whether a given transformation admits a factorization of length at most $l$.
This problem is \PSPACE-complete for~$l$ in binary encoding
and \NP-complete for~$l$ in unary encoding, even for permutation
groups as shown by Jerrum~\cite[Thm.~2.2]{DBLP:journals/tcs/Jerrum85} 
and Even and Goldreich~\cite[Thm.~1]{DBLP:journals/jal/EvenG81}, respectively.
Furthermore, Fernau and Bruchertseifer~\cite{DBLP:journals/mscs/FernauB22}
proved \WSYNC-completeness w.r.t.~$l$ and general monoids.
When asking for a product of length exactly $l$,
Cai et al.~\cite[Thm.~6 \& 7]{DBLP:journals/aml/CaiCDF97} proved~$\textsf{W}[2]$-hardness
in general and $\textsf{W}[1]$-hardness for permutation groups w.r.t.~$l$.
In case of permutation groups, asking
for a product of length exactly $l$
was proved  \DP-complete for a unary encoding of $l$
and \PSPACE-complete
for a binary encoding by Lohrey and Rosowski~\cite[Thm.~4.2 \& 4.3]{DBLP:conf/icalp/LohreyR23}.

For further information on \DFAINT and {\sc Membership},
the paper by Blondin et al.~\cite{BlondinKM16} and the master thesis by Blondin~\cite{Blondin2012msthesis}
give a good overview. For more algorithmic investigations on \DFAINT,
we refer to further work by Wehar and coauthors~\cite{DBLP:journals/corr/abs-2108-05244,OliveiraWehar2018,DBLP:conf/dlt/OliveiraW20,Wehar2014hardness,Wehar2016} and others~\cite{Arrighi2021,DBLP:journals/jcss/FernauHV15,FerKre2017,DBLP:journals/tcs/KarakostasLV03,DBLP:conf/mfcs/LangeR92,DBLP:conf/wia/Wareham00}.
Intersection problems on transformation semigroups 
were investigated by Fleischer and Kufleitner~\cite{DBLP:conf/stacs/FleischerK18} and Fleischer~\cite{DBLP:journals/ijfcs/Fleischer20},
see also the thesis by Fleischer~\cite{Fleischer2019thesis}.
The membership problem for finite monoids and groups 
has been generalized to the \emph{rational (subset) membership problem}~\cite{KAMBITES2007622,Lohrey2015}.



\paragraph{A Remark on Local Confluence and Synchronization} Simon~\cite{Simon1972DotDepthOne,DBLP:conf/automata/Simon75} introduced
piecewise testable languages and characterized them as those languages
acceptable by automata with $\mathcal J$-trivial transition monoids.
Kl{\'{\i}}ma~\cite[Thm.~1]{DBLP:conf/dlt/KlimaP13}
proved that a language is piecewise testable
if and only if it is acceptable by a locally confluent partially ordered automaton.
Since we mentioned locally confluent automata above, 
we point out a connection of synchronizing automata to piecewise testable languages 
that has not been noted before and is based on the following statement.

\begin{proposition}\label{prelim:pto:locally-confluent}
	Let $\automat{A}{Q}{\Sigma}{\delta}{q_0}{F}$ be partially ordered and accessible.
	If $\mathcal A$ is locally confluent, then it is synchronizing.
\end{proposition}
\begin{proof}
	We do induction on $|Q|$.
	Clearly, the claim holds for $1 \le |Q| \le 2$.
	Next, suppose $|Q| \ge 3$.	
	
	Since $\mathcal A$ is accessible,
	there either exists $a \in \Sigma$
	such that $\delta(q_0, a) \ne q_0$
	and $\delta(q_0, a) = \delta(q_0, b)$
	for each $b \in \Sigma$ with $\delta(q_0, b) \ne q_0$,
	or there exist $a,b \in \Sigma$	such that for $p = \delta(q_0, a)$
	and $q = \delta(q_0, b)$ we have $q_0 \ne p$ and $q_0 \ne q$
	and $p \ne q$. In the first case,
	the set $\upperset{\{\delta(q_0, a)\}}$
	is a proper upper set of size at least two.
	In the second case,	by assumption, 
	there exists $u \in \{a,b\}^*$
	with $\delta(p, u) = \delta(q, u)$.
	Since $p \ne q$, we have $p \ne \delta(p, u)$ or $q \ne \delta(q, u)$. 
	Thus either $\upperset{\{p\}}$ or $\upperset{\{q\}}$ has size at least two and neither contains $q_0$.
	
	Thus, in each case we find $q \in Q$ such that  for $S = \upperset{\{q\}}$
	we have $2 \le |S| < |Q|$.\footnote{Alternatively,
		either there exists $a \in \Sigma$
		such that $q = \delta(q_0, a)$ is not a trap state
		and $q \ne q_0$, in which case $\upperset{\{q\}}$
		is proper and has size at least two,
		or $\delta(q_0, a)$ is a trap state for each $a \in \Sigma$.
		Since for each $a,b \in \Sigma$, $\delta(q_0, au) = \delta(q_0, bu)$
		for some $u \in \{a,b\}^*$,
		in the latter case we must have $|Q| \le 2$.}    	
	Then~$\mathcal A_{|S}$ and~$\mathcal A / {\sim_S}$
	are accessible, partially ordered, and locally confluent.
	So, inductively, both automata are synchronizing.
	By \crefitem{prelim:prop:sync_words_restriction}{prelim:Synuv}
	it follows that $\mathcal A$ is synchronizing.
\end{proof}

\begin{corollary}\label{prelim:cor:pto:locally-confluent}
	An accessible partially ordered automaton over an alphabet of size at most two
	is synchronizing if and only if it is locally confluent.
\end{corollary} 
\begin{proof}
	First, let $\semiautomat{A}{Q}{\Sigma}{\delta}$
	be synchronizing with $|\Sigma| \le 2$
	and synchronizing word $u \in \Sigma^*$.
	Let $q \in Q$ and $a,b \in \Sigma$.
	We can assume $a \ne b$,
	so that $\Sigma = \{a,b\}$.
	Thus $u \in \{a,b\}^*$
	and since~$u$ is synchronizing we have $\delta(q, au) = \delta(q, bu)$.
	Hence, $\mathcal A$ is locally confluent.
	
	Since by \cref{prelim:pto:locally-confluent} 
	local confluence implies synchronizability,
	we obtain the claim.
\end{proof}

\begin{figure}[tb]
	\centering 	
	\begin{tikzpicture}[node distance = 1cm and 2.5cm, on grid, auto, ->]
		\node (q0) [state, inner sep=3pt, minimum size=5pt, ] {};
		\node (q1) [state, inner sep=3pt, minimum size=5pt, above right = of q0] {};
		\node (q2) [state, inner sep=3pt, minimum size=5pt, below right = of q0] {};
		
		\path (q0) edge node {$a$} (q1) (q1) edge node {$c$} (q2) (q0) edge [below] node {$b,c$} (q2);
	\end{tikzpicture}
	\caption[Non-locally confluent synchronizing partially ordered automaton]{Synchronizing partially ordered automaton that is not locally confluent. Self-loops omitted.} 
	\label{prelim:classes:fig:pto-not-locally-confluent}
\end{figure}

Thus, with the result by Kl{\'{\i}}ma~\cite[Thm.~1]{DBLP:conf/dlt/KlimaP13} mentioned above,
we obtain that a binary language is piecewise testable
if and only if its minimal automaton is a synchronizing partially ordered automaton.

\Cref{prelim:cor:pto:locally-confluent} does not hold for alphabets with more than two letters,
for example, see \cref{prelim:classes:fig:pto-not-locally-confluent}.
